\chapter{Resource donation}
\label{ch:donation}

\section{Donate everything; protect dynamically}

The instinct when pooling somebody's desktop is to hold capacity back --- donate
half the cores, cap memory at a quarter --- so the owner always has room. This
design does the opposite on workers, and the reason is that static reservation
and the contention policy of Chapter~\ref{ch:policy} defend against the same
thing.

\begin{keypoint}
Holding capacity back only idles it in the common case, which is that nobody is
using the machine. The policy would suspend anyway the moment someone arrives.
Reservation buys nothing the policy does not already provide, and costs
throughput continuously.
\end{keypoint}

That applies to workers. The \textbf{entry node is deliberately different},
because it carries the collector, negotiator and schedd: if the schedd cannot
fork shadows or \kn{fsync} its queue, nothing anywhere in the pool runs.

\begin{table}[H]
\centering
\begin{tabularx}{\textwidth}{L{0.16\textwidth} L{0.19\textwidth} L{0.19\textwidth} Y}
\toprule
\textbf{Resource} & \textbf{Worker} & \textbf{Entry node} & \textbf{Why they differ} \\
\midrule
CPU & \textbf{all threads} & reserve $\sim$half & Daemons need guaranteed headroom \\
Memory & \textbf{90\%} offered & $\sim$35\% offered & One shadow per running job, pool-wide \\
Disk & reserved + tiers & reserved + tiers \textbf{+ transfer caps} & SPOOL is the funnel every job crosses \\
Load scaling & \kn{TotalCpus} & \textbf{\kn{DetectedCpus}} & See \S\ref{sec:detectedcpus} \\
Load thresholds & 0.50 / 0.75 & 0.50 / 0.75 & Now identical; see \S\ref{sec:looser} \\
\bottomrule
\end{tabularx}
\caption{Donation policy: worker versus entry node.}
\end{table}

\section{Scale on DetectedCpus, not TotalCpus}
\label{sec:detectedcpus}

\kn{TotalCpus} is what Condor \emph{advertises}, capped by \kn{NUM\_CPUS}, while
\kn{TotalLoadAvg} is measured across the \emph{whole machine}. On a four-core box
advertising two, scaling load thresholds by \kn{TotalCpus} compares a four-core
load against a two-core yardstick: the machine reads as saturated while half of
it sits idle.

Workers where \kn{NUM\_CPUS} is unset are unaffected only by coincidence. The
entry node, which does cap it, is affected directly --- which is where this was
found.

\section{The entry node's thresholds must be looser}
\label{sec:looser}

\kn{condor\_shadow} processes are children of the \textbf{schedd}, not the
startd, and are therefore \textbf{not} counted in \kn{CondorLoadAvg}.

Every file transfer the entry node brokers for the whole pool consequently lands
in \kn{NonCondorLoad} and looks exactly like owner activity. Worker-tight
thresholds would make the entry node close admission --- and flap its own
running jobs --- purely for doing the job it exists to do.

\section{Fractions scale; absolute headroom does not}

Load thresholds are expressed as \kn{TotalCpus} $\times$ fraction, which travels
across machine sizes. The \emph{slack} such a threshold leaves does not.

At a fraction of 0.25, a sixteen-core box has 4.0 of headroom while a four-core
box has 1.0 --- and a desktop with a browser open already idles above that.

\begin{trap}
\textbf{Measured.} A four-core desktop, otherwise idle, reported non-Condor load
of \textbf{1.23} against a 1.0 threshold. It refused every job it was offered
while \kn{condor\_status} showed it healthy and \st{Unclaimed}. Nothing was
broken; the threshold was simply below the machine's floor.
\end{trap}

The values became 0.50 to admit and 0.75 to suspend. Admitting while the owner
uses half the machine is safe, because \kn{JOB\_RENICE\_INCREMENT} means Condor
jobs lose the CPU to the owner instantly. The gap between the two values is
deliberate hysteresis --- equal thresholds make jobs flap.

\section{Disk floors: absolute on the entry node, fractional on workers}

A fraction suits a volume dedicated to Condor. The entry node's \kn{SPOOL},
\kn{EXECUTE}, \file{job\_queue.log} and operating system share one filesystem
that is already around 86\% full of unrelated data. What the OS needs there is a
\textbf{fixed reserve}, not a percentage of a large, mostly-full disk.

\section{SPOOL cannot be protected by the startd policy}

\kn{condor\_submit -spool} writes into \kn{SPOOL} at \textbf{submit} time,
before any matchmaking happens, so \kn{START} never sees it. The startd policy
--- the entire subject of the next chapter --- is structurally incapable of
defending this.

The only controls are schedd-side: \kn{MAX\_TRANSFER\_INPUT\_MB} and
\kn{MAX\_TRANSFER\_OUTPUT\_MB}. Size them against the disk that actually exists.

\begin{trap}
At 2048\,MB per job with 32\,GB free, \textbf{sixteen concurrent submissions
fill the volume} --- and \kn{MAX\_JOBS\_RUNNING} permits 200. Filling SPOOL
stops the schedd writing its queue, which stops the \textbf{entire pool}.
\end{trap}

\section{Advertisement cadence}

\kn{UPDATE\_INTERVAL} defaults to 300 seconds, a value sized for pools of
thousands of machines. At this scale an advertisement every 30 seconds is free,
and the default costs twice over:

\begin{itemize}
  \item \kn{condor\_status} showed a frozen value for two and a half minutes
        while load was driven from 1 to 15.
  \item After a central-manager restart, workers took roughly twice the interval
        to reappear --- one cycle to discover their cached security session was
        dead, another to re-authenticate (\S\ref{sec:cmrestart}).
\end{itemize}

\begin{keypoint}
Lowering \kn{UPDATE\_INTERVAL} does \textbf{not} change how fast policy reacts.
The startd evaluates \kn{START}, \kn{SUSPEND} and \kn{PREEMPT} every
\kn{POLLING\_INTERVAL} --- five seconds --- regardless. The default made the
policy \emph{invisible} and re-registration slow. It did not make the policy slow.
\end{keypoint}

\section{Sandbox hygiene}

The starter removes a job's scratch directory as soon as output transfer
completes, so the normal path already reclaims immediately.

The leak is \textbf{orphans}: crashed starters, evicted jobs, held jobs. Nothing
reaps those until \kn{condor\_preen} runs, which defaults to \textbf{once every
24 hours}. On a design whose binding constraint is disk, that is the actual
failure mode. Set \kn{PREEN\_INTERVAL} to an hour with \kn{PREEN\_ARGS = -r}.
